\section*{Введение}
\addcontentsline{toc}{section}{Введение}

\subsection*{Актуальность темы исследования}
\addcontentsline{toc}{subsection}{Актуальность темы исследования}

Волатильность финансовых активов является одним из ключевых понятий в современных финансах:
она определяет стоимость опционов, влияет на управление рисками портфелей и служит
индикатором неопределённости на рынке.
Особый интерес представляет волатильность рынка драгоценных металлов~--- золота, серебра,
платины и палладия,~--- которые традиционно выступают защитными активами в периоды
геополитической нестабильности и инфляционного давления.

Начиная с 2022 года российский рынок драгоценных металлов претерпел значительные структурные
изменения: введение западных санкций, ограничение доступа к международным площадкам и
существенное изменение режима волатильности делают прогнозирование на основе
исторических закономерностей недостаточным.
В этих условиях особую ценность приобретают методы, способные учитывать
\textit{настроения инвесторов} и \textit{информационные шоки}~--- реакцию рынка на
новостной поток, изменения ключевой ставки, геополитические события.

\subsection*{Постановка проблемы}
\addcontentsline{toc}{subsection}{Постановка проблемы}

Традиционные модели прогнозирования волатильности опираются в основном на
прошлые ценовые данные и плохо учитывают информационный фон рынка. При этом
неочевидно, помогает ли добавление настроений и внимания инвесторов
прогнозировать волатильность российских драгоценных металлов и какие именно
источники полезны. Эту задачу и решает настоящая работа.

Настоящая работа посвящена разработке системы прогнозирования реализованной волатильности
(Realized Volatility, RV) инструментов Московской биржи. Программный комплекс собирает
данные по всем четырём металлам~--- GLDRUB\_TOM (золото), SLVRUB\_TOM (серебро),
PLTRUB\_TOM (платина) и PLDRUB\_TOM (палладий),~--- однако содержательный
эконометрический анализ сосредоточен на золоте и серебре: для платины и палладия
длина доступного ряда на MOEX недостаточна для надёжной оценки в режиме
расширяющегося окна (см.\ раздел~\ref{sec:results}).
В качестве бенчмарка используется классическая HAR-RV модель~\cite{Corsi2009},
которая расширяется за счёт sentiment- и attention-признаков,
извлекаемых из разнородных источников данных.

\subsection*{Цель и задачи проекта}
\addcontentsline{toc}{subsection}{Цель и задачи проекта}

\textbf{Цель работы}~--- разработать программный комплекс для
прогнозирования реализованной волатильности российских драгоценных металлов,
который объединяет эконометрические и ML-модели с многоисточниковым
вектором настроений и внимания инвесторов и предоставляет
веб-интерфейс для интерактивной визуализации прогнозов и sentiment-индексов.

\textbf{Задачи:}
\begin{enumerate}
    \item Сформировать панельный датасет дневных OHLCV-котировок драгоценных металлов
          на MOEX за период 2018--2026 гг.\ и вычислить реализованную волатильность
          оценщиком Garman--Klass (в коде также реализованы альтернативные
          оценщики Parkinson и Rogers--Satchell).
    \item Собрать многоисточниковый sentiment/attention вектор:
          новостной поток (RSS, GDELT), поисковые запросы (Google Trends),
          объёмы торгов MOEX, макроэкономические индикаторы (ЦБ РФ, FRED).
    \item Оценить HAR-RV бенчмарк и его расширения с sentiment-признаками
          методом расширяющегося окна (expanding window OOS).
    \item Реализовать и обучить модели машинного обучения на том же наборе
          признаков~--- градиентный бустинг XGBoost и kNN-прогнозировщик на
          пространстве информационных состояний (по методологии Halousková
          и Lyócsa)~--- и сравнить их с эконометрическим бенчмарком,
          в том числе по тесту Диболда--Мариано.
    \item Разработать веб-приложение для визуализации прогнозов волатильности
          и sentiment-индексов в режиме реального времени.
\end{enumerate}

\textbf{Научная новизна} работы заключается в адаптации подхода
Halousková и Lyócsa~\cite{Halouskova2025}~--- прогнозирования волатильности
через макроэкономический attention~--- к условиям российского рынка драгоценных металлов
с учётом специфики доступных источников данных (русскоязычный новостной поток,
GDELT, данные MOEX).

\textbf{Практическая значимость.}
Главная ценность работы~--- это законченный программный комплекс:
от сбора семи источников данных и расчёта волатильности до сравнения
четырёх моделей и ве2б-визуализации (более 5000 строк кода,
см.\ раздел~\ref{sec:volume}). Кроме того, в работе сделан единый модуль
метрик с QLIKE~\cite{Patton2011} и тестом Диболда--Мариано~\cite{Diebold1995}.
Видно, что добавление sentiment даёт лишь небольшой
и слабо значимый выигрыш для золота, а классическая HAR-log остаётся сильным
бенчмарком. Этот результат сам по себе полезен для риск-менеджмента.

\subsection*{Обоснование выбора инструментов и технологий}
\addcontentsline{toc}{subsection}{Обоснование выбора инструментов и технологий}

Выбор инструментов определялся задачей и доступными данными. Весь комплекс
написан на Python из-за богатой экосистемы для анализа данных и машинного
обучения (pandas, numpy, scikit-learn, xgboost). В качестве бенчмарка взята
модель HAR-RV, а не GARCH: она проста, интерпретируема и хорошо работает на
дневных данных. Для нелинейного сравнения добавлены XGBoost (стандарт для
табличных данных) и kNN на пространстве информационных состояний (по
методологии Halousková и Lyócsa). Веб-интерфейс сделан на лёгком фреймворке
Flask с интерактивными графиками Plotly. Все источники данных~--- бесплатные и
общедоступные (MOEX ISS, GDELT, Google Trends, ЦБ~РФ, FRED, RSS, Telegram),
что обеспечивает воспроизводимость результатов.

